\subsection{One referral after tentative offers}\label{sec:residualreferral}
An applicant sees five programs, applies, and receives no offer. A clearinghouse can now observe which seats are still vacant. Suppose it can also see the applicant's eligible alternatives or ask which additional programs they would attend. How many people could receive a place with one additional referral each and all provisional contracts still feasible?

Let $D$ be the initial display, $M$ any stable matching on $D$, and $U$ the unmatched applicants. Program $j$ has integer capacity $c_j$ and $v_j=c_j-|M(j)|$ vacancies. Its fixed, strict priorities are responsive: given any set of acceptable applicants, it prefers to fill its capacity with the highest-priority members of that set. Construct a graph between $U$ and programs with $v_j>0$, keeping each mutually acceptable pair in the full feasible graph. Each such pair is hidden in $D$; otherwise it would block $M$. Assume $q_i=\min(k,d_i)$, where $k\ge1$ and $d_i$ is the applicant's feasible degree. Give program $j$ capacity $v_j$ in this graph and let $r$ be its maximum matching size.

\begin{theorem}[Vacancy referrals with a fixed quota]\label{thm:residualreferral}
At most $r$ initially unmatched applicants can be added to an allocation containing every original contract. This number is attainable with one display substitution for each of $r$ applicants and no change to any other list. Applicant-proposing DA on the resulting display places exactly the original recipients and those $r$ applicants and gives each original recipient a weakly preferred contract. If only a subset of the selected applicants agrees to its substitution, these same conclusions hold for that subset.
\end{theorem}
\begin{proof}
Any extension keeping $M$ intact can add at most $v_j$ people at each program. Its added pairs therefore form a matching in the residual graph. Conversely, take any residual matching $A$. Show each selected applicant their assigned program, dropping their least preferred displayed program to keep their quota. Under the usual rule $q_i=\min(k,d_i)$ the existence of a hidden feasible pair implies $d_i>k$; if the platform instead allows an unused slot, it can use that slot.

Write $W=M\cup A$. Each newly displayed pair belongs to $W$. On an old displayed pair the applicant's assignment has weakly improved relative to $M$. An originally full program has exactly its old occupants in $W$. At an originally slack program every applicant with an old displayed edge already weakly preferred their assignment under $M$ to that program, by stability. Thus no displayed pair blocks $W$, and all partners remain visible. Applicant-proposing DA weakly dominates $W$ for every applicant, and invariance of the matched applicant set across stable outcomes gives exactly its recipient set. For opt-outs delete the relevant edges from $A$ and keep their original menus; the proof applies to the remaining residual matching.
\end{proof}

This is a certificate for an extension that leaves the original contracts available as a stable witness. DA may improve an incumbent's partner. Moving incumbents on their own newly revealed edges could create further placements, an objective outside this bound. A prior year's predicted vacancies do not certify an extension of today's offers. The stability guarantee is relative to the revised display. The institution needs the ability to direct individual referrals or to reserve them, alongside an accepted list of otherwise undisplayed choices. Publicly announcing every vacancy can induce several people to chase the same seat. \citet{lepesure} simulates a Chilean school mechanism that helps unassigned applicants find previously unlisted nearby schools and runs a targeted residual submarket preserving earlier allocations. Our question concerns previously submitted options displaced by a quota. We compute the maximum contract extension under that restriction and certify a full-market DA rerun after an individual substitution or opt-out. Other work studies efficiency after cancellations and repeated general admission rounds \citep{feigenbaumreallocation,gradualschool}.

The graph also determines whether access has to be rationed among people with a feasible vacancy. For $L$ a subset of unmatched applicants with a residual edge and $N(L)$ their neighboring programs, at most $\sum_{j\in N(L)}v_j$ members of $L$ receive one of those seats. The largest deficit $|L|-\sum_{j\in N(L)}v_j$ is exactly the number of applicants with a residual edge left out by a maximum flow. A minimum cut provides a witness. If a group of applicants is to receive priority among maximum-size extensions, solve a maximum-weight residual matching with weight $n+1+\mathbf{1}\{i\text{ belongs to the group}\}$ on every eligible edge. The integer $n$ bounds the number of applicants. This gives the most group members who can be served at the maximum total; the same witness proof applies.
